
\section{Discussion, Limitations, and Future Directions}
\label{sec:discussion}

\subsection{The Teacher-Choice Effect}

The ablation study showcases a clear rank ordering under Protocol~B (where $\succ$ denotes strictly lower mean total queue length): Reflected UAS teacher ($9.82$) $\succ$ UAS teacher ($11.14$) $\succ$ scratch ($5{,}451$; diverged).
The $1.32$-unit gap between the Reflected UAS teacher and the UAS teacher exceeds the impact of the three tested architectural modifications (log-normalization, zero-init, scratch initialization), which identifies teacher quality as a primary design factor in this empirical setting.

This observation has a natural interpretation.
Essentially, behavior cloning positions the neural policy's parameters in a region from which \reinforce{} fine-tuning can make only modest improvements under the tested budget.
A competent teacher provides a better initialization, from which \reinforce{} then attempts to drive further improvement, whereas a weaker initialization forces fine-tuning to expend its efforts on overcoming the deficit.
Within the tested optimization budget, the learned policy behaves more like local refinement of the initialization than like global policy search.
The offline BC phase does not address the distributional shift~\cite{ross2011reduction}.
On-policy data collection during pretraining could further improve initialization quality.

\subsection{Why Reflected UAS Outperforms UAS}

Reflected UAS alters three key components of UAS:
\begin{itemize}
\item \textbf{Tempered normalization} ($\calbeta = 0.85$): the denominator $\srvrate_i^{\calbeta}$ increases more slowly than $\srvrate_i$, reducing the routing-mass differential between fast and slow servers.
\item \textbf{Softened prior} ($\calgamma = 0.5$): the pre-exponential factor $\srvrate_i^{\calgamma}$ reduces the routing bias toward fast servers relative to UAS's proportional-to-$\srvrate_i$ prior.
\item \textbf{Reduced look-ahead} ($\caloffset = 0.5$): the energy function's additive constant is decreased from $1$ (UAS: $\Qlen{i}+1$) to $0.5$, shifting routing slightly toward current-state responsiveness.
\end{itemize}
These mechanisms are consistent with the empirical improvement---the tempered normalization and softened prior jointly reduce over-routing to fast servers (Gini $0.174$ vs.
UAS's $0.328$), while the diminished offset increases responsiveness to current state---but a formal attribution of the $1.45$-unit improvement ($10.042$ vs.
$11.495$) to individual parameters remains open.

\subsection{Limitations}

\begin{enumerate}[label=\textbf{L\arabic*.}, leftmargin=2.2em]
\item \textbf{Full Reflected UAS family remains open:}
Section~\ref{sec:reflected_uas} identifies a theorem-backed certified subset within the reflected family under the explicit sufficient condition $\varepsilon_{\calbeta,\calgamma,\caloffset}>0$.
However, our audit in Table~\ref{tab:sruas_audit} finds that the benchmark default point $(\calbeta, \calgamma, \caloffset) = (0.85, 0.5, 0.5)$ yields $\varepsilon_{\calbeta,\calgamma,\caloffset} < 0$ under the anchor configuration.
While we provide \textbf{numerical evidence} that a benchmark-specific piecewise candidate has negative drift on the verified boundary shells (Table~\ref{tab:piecewise_candidate}), a formal symbolic proof for the unrestricted family remains an open theoretical question. This establishes a definitive split in our results: the benchmark-default policy is a competitive empirical baseline, but it is not yet theorem-certified.

\item \textbf{Default learned configuration is not the strongest:}
The standard $\ngibbsq$ model (trained with a UAS teacher) achieves $11.68$ in the anchor benchmark, underperforming JSSQ ($11.02$) and Reflected UAS ($10.04$).
The configuration with the reflected teacher ($9.82$) was identified through ablation.
These findings confirm that the choice of teacher policy is the dominant factor in policy-gradient performance, but they also clarify that our reported neural "best" refers specifically to this reflected-teacher variant, which requires separate behavior-cloning initialization. Additionally, interaction effects between ablation factors were not exhaustively tested; for instance, the combination of the Reflected UAS teacher with architectural modifications (e.g., removing Log-Normalization) may yield further improvement.
Crucially, absolute values from Protocol A and the ablation study (Protocol B) reflect different simulation horizons ($15{,}000$ vs.\ $20{,}000$ units) and are not directly comparable.

\item \textbf{Competitive comparisons limited to $\Nservers = 10$:}
The stress test (Appendix~\ref{sec:appendix_stress}) validates UAS scalability up to $\Nservers = 1{,}024$, however neither Reflected UAS nor $\ngibbsq$ are evaluated competitively at $\Nservers > 10$. The primary constraint is training stability and computational cost: the \reinforce{} policy gradient estimates become increasingly noisy as the action space and state dimension grow with $\Nservers$, and the behavior-cloning phase requires more teacher samples. Investigating architecture choices (e.g., attention over server states, equivariance) that scale to larger $\Nservers$ is an important open problem.
We scope the comparative evaluation of the empirical routing pipeline to $\Nservers = 10$. This highlights a deliberate split in our evaluation: while the $\ngibbsq$ empirical pipeline performs strongly at the benchmark scale, no scaling guarantee is claimed for $\ngibbsq$ or Reflected UAS beyond the evaluated $\Nservers=10$. The analytical UAS policy, by contrast, is separately stress-tested up to $\Nservers = 1{,}024$ without neural retraining.

\item \textbf{Coarse parameter grid for Reflected UAS:}
The benchmark parameters $(\calbeta, \calgamma, \caloffset) = (0.85, 0.5, 0.5)$ were selected from a grid of 64 configurations evaluated on a held-out validation set.
The grid is coarse and the parameter surface may not be convex.
Therefore, a more systematic search method (e.g., Bayesian optimization) could enhance the confidence in the improvements observed.

\item \textbf{REINFORCE gradient magnitude error:}
The \reinforce{} estimator demonstrates 32.6\% relative magnitude error versus finite differences (Appendix~\ref{sec:appendix_gradient}).
While directionally validated (cosine similarity $1.0000$), the gradient signal remains somewhat noisy. The magnitude mismatch effectively reduces the learning rate, demonstrating that the full \reinforce{} training story is not fully solved and variance reduction techniques remain necessary.
Variance reduction techniques such as generalized advantage estimation (GAE) or actor-critic baselines could reduce this error; trust-region methods such as PPO would additionally stabilize step sizes.

\item \textbf{Best-of-5 seed selection bias:}
The strongest neural result ($9.82$ under Protocol B) is the best of 5 random initialization seeds selected on a held-out validation set.
This selection introduces optimistic bias whose magnitude cannot be quantified without per-seed logs (which were not preserved).
The unbiased ensemble mean remains unknown.

\item \textbf{Near-critical load experiments underpowered:}
The near-critical load comparison (Table~\ref{tab:critical}) uses only $n{=}2$ replications per load factor, which is insufficient for statistically reliable confidence intervals.
The reported improvement percentages (8.8\%--22.2\%) are indicative directional evidence, not conclusive quantitative claims.

\item \textbf{No independent seed-range replication:}
All experiments use seeds 42--73. A sensitivity check across an independent seed block (e.g., 100--131) was not performed, so all results depend on a single seed range.

\item \textbf{Markovian assumptions:}
All stability proofs and empirical evaluations assume Poisson arrivals and exponential service times. Robustness to non-Poisson traffic or heavy-tailed service distributions remains unvalidated.

\item \textbf{No comparison to standard RL algorithms:}
The ablation study compares only within the REINFORCE family. Whether more sophisticated RL algorithms (e.g., PPO, SAC) could achieve stronger results from the same initialization remains an open question. Our design isolates the teacher-effect variable but does not establish that REINFORCE is the optimal RL algorithm for this setting.

\item \textbf{Reflected UAS parameters are benchmark-specific:}
The parameters $(\calbeta, \calgamma, \caloffset) = (0.85, 0.5, 0.5)$ were tuned on a single benchmark configuration. Transfer to other server configurations, load regimes, or arrival patterns is not guaranteed.

\item \textbf{Neural policy lacks formal stability guarantee:}
The $\ngibbsq$ policy, being a neural network, has no formal stability certification. It could exhibit instability under distribution shift or in operating regimes far from the training distribution.
\end{enumerate}

\paragraph{When GibbsQ policies may perform poorly}
Entropy-regularized routing may degrade when: (a)~servers are near-homogeneous, where softmax exploration is wasteful and JSQ is near-optimal; (b)~load is very low ($\rho < 0.5$), where policy choice has minimal impact; or (c)~load is extremely high ($\rho > 0.99$), where all soft policies approach instability and greedy policies like JSQ are preferable. The temperature parameter trades off stability (guaranteed for all $\temp > 0$) against efficiency (improved as $\temp \to 0$ for near-homogeneous systems).

\subsection{Future Directions}

\begin{enumerate}[label=\textbf{F\arabic*.}, leftmargin=2.2em]
\item \textbf{Beyond the certified subset:}
Extending the certified-subset Lyapunov argument to cover the full three-parameter family $(\calbeta, \calgamma, \caloffset)$ would close the remaining gap between the analytical and empirical layers.
The core challenge is to enlarge the theorem-backed region beyond the explicit sufficient condition $\varepsilon_{\calbeta,\calgamma,\caloffset}>0$ without sacrificing a tractable drift identity, ideally far enough to certify practically competitive reflected points under the anchor benchmark.
The piecewise candidate in \cref{eq:piecewise_candidate} suggests that a boundary-refined Foster--Lyapunov construction based on quadratic, cubic, and imbalance terms is a concrete route for that extension.

\item \textbf{Competitive comparisons at intermediate scale:}
Competitive policy comparisons at $\Nservers \in \{32, 64, 128\}$---which connects the primary benchmark ($\Nservers = 10$) to the stress test ($\Nservers = 128{+}$)---would establish whether the learned-policy advantage persists as server count grows.

\item \textbf{Advanced RL training:}
Variance reduction via GAE or actor-critic baselines would address the \reinforce{} magnitude error identified in L5.
However, trust-region methods (PPO, natural policy gradient) could additionally improve step-size stability and accelerate convergence.
An immediate empirical question is whether stronger RL optimizers preserve the same teacher-choice ordering observed here, or materially reduce the gap between UAS-teacher and reflected-teacher initialization.

\item \textbf{Non-stationarity:}
Fluctuating arrival rates and dynamic server availability would test whether the GibbsQ stability guarantees and routing policies extend to practical non-stationary settings. A practical online adaptation layer---e.g., lightweight continual fine-tuning of the $\ngibbsq$ network as statistics shift---would bridge the gap to deployed adaptive routing, though such adaptation must be balanced against stability risk during the adaptation transient.

\item \textbf{Sharper mixing-rate results:}
For UAS, the current manuscript establishes positive Harris recurrence with explicit drift constants, but not a geometric-drift theorem. A useful next step would be to derive sharper convergence-rate bounds, and to determine whether geometric ergodicity can be proved for UAS or extended to broader reflected-policy classes.

\item \textbf{Heavy-traffic and diffusion limits:}
A heavy-traffic diffusion-limit analysis (e.g., Halfin--Whitt scaling) would substantially strengthen the theoretical contribution by characterizing asymptotic performance as $\rho \to 1$, but is outside the current scope.

\item \textbf{Critical validation experiments:}
Three experiments would materially strengthen the evidence stack: (a)~a PPO-from-scratch training run, which would determine whether the observed scratch-divergence behavior is REINFORCE-specific or more fundamental; (b)~a non-Markovian robustness test (e.g., Erlang-2 arrivals), which would test whether the policy ranking is preserved under mild departures from Poisson/exponential assumptions; and (c)~an independent seed block (e.g., seeds 1000--1031) for the primary comparison, which would rule out seed-range dependence.
\end{enumerate}
